\begin{abstract}
同策略蒸馏（OPD）通过监督学生自身策略采样的轨迹来传递教师的能力，然而其泛化行为仍不明确，因为大多数研究仅在单一领域和与训练数据相近的基准上评估OPD。我们提出一项受控研究，逐个变化一个泛化因子，从域内分布偏移到跨领域迁移以及多教师设置。我们发现，OPD传递的是教师的推理行为而非对特定问题的答案：训练难度几乎无关紧要，甚至教师从未解决过的问题也具有价值。迁移效果强烈依赖于教师与学生之间的原点关系：同源配对能使学生在不同语言、推理范围甚至其他领域中接近教师，而异源配对则主要拟合训练分布。

这种广泛的影响是一把双刃剑：由于将提示路由给领域专家无法限制每位教师的影响范围，因此组合它们会带来能力之间的混合依赖性\emph{跷跷板}效应。这些结果明确了OPD何时能够泛化，并为诊断多教师OPD提供了有用的视角。
\end{abstract}

